\subsection{有界测度的窄收敛}

设 \(\mathbf{T}\) 是完全正则空间；双线性形式

\[
(f, \mu) \mapsto \int f (t) d \mu (t)
\]

在 \(\mathcal{C}^{b}(\mathrm{T})\times\mathcal{M}^{b}(\mathrm{T})\) 上使这两个空间构成分离对偶。事实上，该对偶在 \(\mathcal{C}^{b}(\mathrm{T})\) 中显然是分离的，因为测度 \(\varepsilon_{x}\;(x\in\mathrm{T})\) 属于 \(\mathcal{M}^{b}(\mathrm{T})\)；它在 \(\mathcal{M}^{b}(\mathrm{T})\) 中也是分离的（第 1 号命题 2）。

定义 1. --- \(\mathcal{M}^{b}(\mathrm{T})\) 上与前述 \(\mathcal{C}^{b}(\mathrm{T})\) 和 \(\mathcal{M}^{b}(\mathrm{T})\) 之间的对偶相联系的弱拓扑，称为 \(\mathcal{M}^{b}(\mathrm{T})\) 上的窄收敛拓扑（或窄拓扑）。

窄拓扑是豪斯多夫的，这是由定义前的备注得到的。我们将经常使用副词“按窄拓扑地”，意指“按窄拓扑的意义”。除非另有说明，在本节余下部分始终给 \(\mathcal{M}^b (\mathrm{T})\) 赋予窄拓扑。

\(\mathcal{C}^{b}(\mathrm{T})\) 的每个元素都是 \(\mathcal{C}_{+}^{b}(\mathrm{T})\) 中元素的线性组合。要使 \(\mathcal{M}^{b}(\mathrm{T})\) 上的滤子 F 窄收敛到有界测度 \(\lambda\)，必要且充分的条件是

\begin{enumerate}
\def\labelenumi{(\arabic{enumi})}
\setcounter{enumi}{3}
\tightlist
\item
\end{enumerate}

\(\lim_{\mu} \mu(f) = \lambda(f)\) 关于 \(\mathfrak{F}\) 成立，对每个 \(f \in \mathcal{C}_+^b(\mathbf{T})\) 都成立。

备注。 --- 1) 若 T 局部紧，则窄拓扑比由淡拓扑在 \(\mathcal{M}^{b}(\mathrm{T})\) 上诱导的拓扑细，并且这两个拓扑仅当 T 紧时才一致。事实上，若 T 不紧，映射 \(t \mapsto \varepsilon_{t}\) 关于 T 的相对紧子集的补集所成的滤子淡收敛到 0，但不按窄拓扑收敛到 0，因为函数 1 属于 \(\mathcal{C}^{b}(\mathrm{T})\)（关于淡收敛与窄收敛的关系，参见命题 9）。

\begin{enumerate}
\def\labelenumi{\arabic{enumi})}
\setcounter{enumi}{1}
\item
  由命题 4 立即得到 \(\mathcal{M}_+^b (\mathbf{T})\) 在 \(\mathcal{M}^b (\mathbf{T})\) 中是闭的。
\item
  若 \(\mathbf{T}\) 完全正则，则映射 \(t \mapsto \varepsilon_t\) 从 \(\mathbf{T}\) 到 \(\mathcal{M}^b(\mathbf{T})\) 是一个同胚（GT, IX, §1, No.~5）。
\end{enumerate}

命题 6. --- 设 T 是完全正则空间。

\begin{enumerate}
\def\labelenumi{\alph{enumi})}
\item
  设 \(f\) 是满足 \(\geqslant 0\) 且定义在 \(\mathbf{T}\) 上的下半连续数值函数；则映射 \(\mu \mapsto |\mu|^{\bullet}(f)\) 在 \(\mathcal{M}^b(\mathbf{T})\) 上下半连续。
\item
  设 \(f\) 是定义在 \(\mathbf{T}\) 上的上半连续有界函数；则映射 \(\mu \mapsto \mu(f)\) 在 \(\mathcal{M}_+^b(\mathbf{T})\) 上上半连续。
\end{enumerate}

事实上，由第 1 号命题 b)，\(\mu \mapsto |\mu|^{\bullet}(f)\) 是形如 \(\mu \mapsto |\mu(g)|\)（其中 \(g \in \mathcal{C}^b(T)\)）的函数族的上包络，因而关于窄拓扑连续。这就证明了 \(a\)。为证明 \(b\)，只须取 \(f\) 的一个常数上界 C，并写成 \(\mu(f) = \mu(C) - \mu(C - f)\)；映射 \(\mu \mapsto \mu(C)\) 连续，而根据上文，映射 \(\mu \mapsto \mu(C - f)\) 在 \(\mathcal{M}_+^b(T)\) 上下半连续。

命题 7. --- 设 T 是完全正则空间。设 \(\mu\) 是 T 上的有界正测度，f 是 T 上的有界正函数，并且 T 中 f 不连续的点集局部地对 \(\mu\) 可忽略。则映射 \(\lambda \mapsto \lambda^{\bullet}(f)\) 从 \(\mathcal{M}_{+}^{b}(\mathrm{T})\) 到 R 在点 \(\mu\) 处连续。

对每个 \(t \in \mathbb{T}\)，置 \(f'(t) = \liminf_{s \to t} f(s)\)、\(f''(t) = \limsup_{s \to t} f(s)\)。显然 \(f' \leqslant f \leqslant f''\)，并且在 \(\mathbb{T}\) 中 \(f\) 连续的每一点处等号成立（因而 \(\mu\)-几乎处处成立）。另一方面，\(f'\) 下半连续，\(f''\) 上半连续且有界（GT, IV, §6, No.~2, 命题 4）。由命题 6，因而有以下关系，

\[
\begin{array}{c} \mu^ {\bullet} (f ^ {\prime}) \leqslant \liminf _ {\lambda \to \mu} \lambda^ {\bullet} (f ^ {\prime}) \leqslant \liminf _ {\lambda \to \mu} \lambda^ {\bullet} (f) \leqslant \mu^ {\bullet} (f) \leqslant \limsup _ {\lambda \to \mu} \lambda^ {\bullet} (f) \\ \leqslant \limsup _ {\lambda \to \mu} \lambda^ {\bullet} (f ^ {\prime \prime}) \leqslant \mu^ {\bullet} (f ^ {\prime \prime}). \end{array}
\]

最后注意到 \(\mu^{\bullet}(f') = \mu^{\bullet}(f'')\)，因为 \(f'\) 与 \(f''\) 局部地在 \(\mu\)-几乎处处相等。

命题 8. --- 设 X 是完全正则空间，T 是 X 的子空间，i 是从 T 到 X 的典范嵌入。记 W 为集中于 T 的 X 上有界正测度所成的集合，并赋予它由 \(\mathcal{M}^{b}(X)\) 诱导的拓扑。则映射 \(\mu \mapsto i(\mu)\) 从 \(\mathcal{M}_{+}^{b}(T)\) 到 \(\mathcal{M}^{b}(X)\) 是 \(\mathcal{M}_{+}^{b}(T)\) 到 W 的同胚。

我们仍以 \(i\) 表示映射 \(\mu \mapsto i(\mu)\)，它从 \(\mathcal{M}_{+}^{b}(\mathrm{T})\) 到 \(\mathcal{M}_{+}^{b}(\mathrm{X})\)；\(i\) 是单射（§2, No.~4, 命题 8），并将 \(\mathcal{M}_{+}^{b}(\mathrm{T})\) 映入 W（§2, No.~3, 命题 7）。若 \(\lambda \in \mathrm{W}\)，则 \(\lambda = i(\lambda_{\mathrm{T}})\)（§2, No.~3, 命题 7 b)）。因此，\(i\) 是 \(\mathcal{M}_{+}^{b}(\mathrm{T})\) 到 W 的双射，而 \(i\) 的逆双射是 W 上的映射 \(r : \lambda \mapsto \lambda_{T}\)。另一方面，i 连续：事实上，若 \(f \in \mathcal{C}^{b}(X)\)，则 \(\langle i(\mu), f \rangle = \langle \mu, f \circ i \rangle\)，且 \(f \circ i\) 属于 \(\mathcal{C}^{b}(T)\)。所以归结为证明：对每个测度 \(\mu \in W\) 及每个函数 \(f \in \mathcal{C}_{+}^{b}(T)\)，有

\[
\lim _ {\lambda \to \mu , \lambda \in W} \lambda_ {T} (f) = \mu_ {T} (f),
\]

或者写成

\[
\lim _ {\lambda \to \mu ,   \lambda \in \mathrm{W}} \lambda (f ^ {0}) = \mu (f ^ {0}).
\]

令 \(f^{\infty}\) 为 X 上的函数，它在 T 上与 f 一致，在 X - T 上等于 \(+\infty\)；令 \(f'\) 和 \(f''\) 分别为 \(f^{0}\) 的上半连续正则化和 \(f^{\infty}\) 的下半连续正则化（GT, IV, §6, No.~2）。关系

\[
f ^ {\prime} (x) = \limsup _ {y \to x} f ^ {0} (y), \quad f ^ {\prime \prime} (x) = \liminf _ {y \to x} f ^ {\infty} (y)
\]

立即推出 \(f'\) 和 \(f''\) 在 T 上都分别与 f 和 \(f^{0}\) 一致。于是命题 6 给出

\[
\mu^ {\bullet} (f ^ {\prime}) \geqslant \lim _ {\lambda \to \mu ,   \lambda \in W} \sup \lambda^ {\bullet} (f ^ {\prime})  , \quad \mu^ {\bullet} (f ^ {\prime \prime}) \leqslant \lim _ {\lambda \to \mu ,   \lambda \in W} \inf \lambda^ {\bullet} (f ^ {\prime \prime})  .
\]

但是，可以在这两个公式中用 \(f'\) 和 \(f''\) 替换为 \(f^{0}\)，因为测度 \(\lambda\) 和 \(\mu\) 集中于 T；于是得到所需关系。

命题 8 的陈述仅对正测度有效：映射 \(\mu \mapsto i(\mu)\) 从 \(\mathcal{M}^b (\mathrm{T})\) 到 \(\mathcal{M}^b (\mathrm{X})\) 是单射且连续的，但一般不是 \(\mathcal{M}^b (\mathrm{T})\) 到其像的同胚。例如，取 \(\mathbf{X} = \mathbf{R}\)、\(\mathrm{T} = \mathbf{R} - \{0\}\)；测度 \(\lambda_t = \varepsilon_t - \varepsilon_{-t}\)（\(t > 0\)）当 \(t\) 趋于 0 时在 X 中按窄拓扑收敛到 0，但在 T 中不按窄拓扑收敛到 0（集合 {]}0, +∞{[} 的特征函数属于 \(\mathcal{C}^b (\mathrm{T})\)）（不过参见第 5 号定理的推论）。

命题 9. --- 设 T 是局部紧空间，F 是 \(\mathcal{M}_{+}^{b}(T)\) 上淡收敛到有界测度 \(\mu\) 的滤子。F 窄收敛到 \(\mu\) 的必要且充分条件是：\(\lim_{\lambda}\lambda(1)=\mu(1)\) 关于 F 成立。

必要性显然。为证明充分性，记 X 为 T 的 Alexandroff 紧化（GT, I, §9, No.~8），i 为从 T 到 X 的典范嵌入。由命题 8，归结为证明映射 \(\lambda \mapsto i(\lambda)\) 关于 F 窄收敛到 \(i(\mu)\)（在 \(\mathcal{M}^{b}(X)\) 中）。由于 \(\mu(1) < +\infty\)，存在 \(A \in F\)，使得 A 中各测度的总质量都被某个数 M 控制；因此只须验证

\[
\lim _ {\lambda , \mathfrak {F}} \int_ {\mathrm{X}} g d (i (\lambda)) = \int_ {\mathrm{X}} g d (i (\mu))\tag{5}
\]

对函数 \(g \in \mathcal{C}^b(\mathrm{X})\)（它们在 \(\mathcal{C}^b(\mathrm{X})\) 中构成总集），上述等式成立。现在，当 \(g\) 在 \(\mathrm{T}\) 中具有紧支集时，由于 \(\mathfrak{F}\) 淡收敛到 \(\mu\)，上述等式成立；当 \(g\) 是 \(\mathrm{X}\) 上的常函数时，由于 \(\lim_{\lambda, \mathfrak{F}} \lambda(1) = \mu(1)\)，上述等式也成立。由于前述两类函数在 \(\mathcal{C}^b(\mathrm{X})\) 中构成总集（第 III 章 §1, No.~2, 命题 3），证明完毕。

